\section{Causal Actionable Counterfactuals}
\label{sec:method}

\cac{} makes causal validity structural rather than optional. It never names a
destination point; it searches for an \emph{action}, and it scores that action on
the record its causal model says the action produces. The audited object and the
optimised object are therefore the same record $x^{a}$ of
Eq.~\eqref{eq:propagate}, so a plan that would not survive its own consequences is
not a low-scoring candidate --- it is not the plan \cac{} is scoring at all.

\subsection{Objective}
A candidate action $a$ acts on a feature subset $S$ with $|S|\le K$
($K=\text{MAX\_ACTIONS}=3$) and is scored through the resulting record
$x^{a}=\mathcal{M}(x^{0};\dopart{X_S=v_S})$ by
\begin{equation}
  J(a) \;=\; \lambda_{\text{v}}\,\underbrace{\big[\tau+\delta-f(x^{a})\big]_{+}}_{\text{validity}}
        \;+\; \lambda_{\text{c}}\,\underbrace{c(a)}_{\text{cost}}
        \;+\; \lambda_{\text{f}}\,\underbrace{\phi(a)}_{\text{feasibility}}
        \;+\; \lambda_{\text{p}}\,\underbrace{\rho(x^{a})}_{\text{plausibility}},
  \label{eq:objective}
\end{equation}
minimised over $a$, with $[\,\cdot\,]_{+}=\max(0,\cdot)$ and weights
$(\lambda_{\text{v}},\lambda_{\text{c}},\lambda_{\text{f}},\lambda_{\text{p}})
=(1.0,\,0.6,\,0.3,\,0.2)$.
\begin{itemize}
  \item \textbf{Validity} is a hinge on the target score that vanishes once
    $f(x^{a})$ clears the threshold $\tau$ by a margin
    $\delta=\text{CAC\_TARGET\_MARGIN}=0.05$. Because it is evaluated on the
    propagated record, driving this term to zero \emph{is} causal validity ---
    there is no separate ``flip the named point'' to fall short of.
  \item \textbf{Cost} $c(a)\in[0,1]$ is the normalised effort of the actions
    (\S\ref{sec:cost}), priced on what the person does, not on the consequences
    the model adds.
  \item \textbf{Feasibility} $\phi(a)$ adds the \texttt{INFINITE\_COST}
    admissibility guard (an inadmissible action is unbounded here, so it can never
    win) and a penalty for disturbing sensitive attributes downstream.
  \item \textbf{Plausibility} $\rho(x^{a})$ is the mean distance from $x^{a}$ to
    its $k{=}5$ nearest favourable-class neighbours, normalised by the dataset
    median, so the produced record resembles a real accepted applicant.
\end{itemize}

\subsection{Search}
The objective is not differentiable and cannot be made so: the black box is opaque
by construction (\S\ref{sec:problem}), the choice of which features to act on is a
discrete subset selection, and the causal model is piecewise-constant wherever a
feature is categorical. \cac{} therefore uses a gradient-free elitist population
search (Algorithm~\ref{alg:cac}), in the elitist tradition of
\citet{deb2002nsga}. A candidate is a mask marking the acted-on
features together with their target values; the population is evaluated in one
batch through the causal model, the better half is retained, and the rest are
regenerated by crossover, Gaussian jitter on the values, and a set mutation that
adds or drops a feature from the acting set. Every child is \emph{legalised} ---
clipped to admissible bounds and grids, stripped of immutable features, and capped
at $K$ actions --- so the search only ever scores plans a person could take. The
$k$ best distinct survivors are returned as the applicant's ranked suggestions.

\begin{algorithm}[t]
\caption{\cac{} action search for one rejected applicant.}
\label{alg:cac}
\begin{algorithmic}[1]
\Require black box $f$, causal model $\mathcal{M}$, applicant $x^{0}$, target class,
  cost model, favourable data
\State $n \gets \max(\text{CAC\_POPULATION},\,\text{MIN\_POPULATION}) = 60$;\quad
  $m \gets \max(2,\lceil n\cdot\text{ELITE\_FRAC}\rceil)=30$
\State initialise $n$ legal candidates: random acting sets of size $\le K$ with
  values drawn from each feature's admissible grid
\For{$g = 1$ \textbf{to} $\text{CAC\_MAX\_ITERS}=200$}
  \For{each candidate $a$}\Comment{one batched pass}
    \State $x^{a} \gets \mathcal{M}(x^{0};\dopart{X_S=v_S})$
      \Comment{abduction--action--prediction}
    \State score $J(a)$ by Eq.~\eqref{eq:objective}
  \EndFor
  \State $E \gets$ the $m$ candidates of lowest $J$\Comment{elitist retention}
  \State children $\gets$ from $E$: uniform crossover, then value jitter
    $\mathcal{N}(0,\text{CAC\_LEARNING\_RATE}{=}0.1)$, then set mutation
    (prob.\ $0.35$)
  \State legalise every child (bounds, grids, mutability, $|S|\le K$)
  \State population $\gets E \cup \text{children}$
\EndFor
\State \Return the $k$ distinct legal candidates of lowest $J$
\end{algorithmic}
\end{algorithm}

\subsection{A note on reconstruction}
The specification we built from was truncated inside the subsection that gave
\cac{}'s pseudocode, so the objective of Eq.~\eqref{eq:objective} and the search of
Algorithm~\ref{alg:cac} are a faithful reconstruction consistent with the parts of
the specification that survived --- the four named terms, their weights, the
action budget, and the intervention-through-SCM scoring. This is stated plainly
among the deviations in the README so no reader mistakes the reconstruction for a
verbatim reproduction. What matters for the claims that follow is the property the
reconstruction preserves exactly: \cac{} scores interventions through the causal
model, so it cannot propose a world its own model disowns.
